% Modified by LianTze Lim to work with fontspec/xelatex
    % Choose for notes or slides
    \if \slidemode 1
        \documentclass{beamer}
        \else
        \documentclass{article}
        \usepackage{beamerarticle}
        \usepackage[a4paper,margin=2cm]{geometry}
        \usepackage{multicol}
    \fi 
\usepackage{amsfonts,amsmath,amssymb,amsthm,amscd}
\usepackage[style=authoryear]{biblatex}
\renewcommand*{\nameyeardelim}{\addcomma\addspace}
\AtBeginBibliography{\tiny}

\usetheme{DarkConsole}
\setbeamertemplate{theorems}[normal font]

% Packages
\usepackage{hyperref}
\usepackage{graphicx}
\usepackage{multimedia} % provides \movie
\usepackage{tikz,pgfplots}
\usetikzlibrary{arrows.meta,calc}
\pgfplotsset{compat = 1.18}
\usepgfplotslibrary{fillbetween}
\usepackage{physics}
\usepackage{siunitx}
\usepackage{polynom}
\usepackage{biblatex}
\usepackage[outline]{contour} % glow around text
\usetikzlibrary{angles,quotes,intersections,fillbetween} % for pic
\newcommand{\myscale}{0.7}
\contourlength{1.3pt}

\tikzset{>=latex} % for LaTeX arrow head
\usepackage{xcolor}
\tikzstyle{vector}=[->,very thick,xcol]
\tikzstyle{mydashed}=[dash pattern=on 2pt off 2pt]
\def\tick#1#2{\draw[thick] (#1) ++ (#2:0.1) --++ (#2-180:0.2)} 

% helper: open-dot macro (hollow circle)
  \newcommand{\opendot}[2]{%
    \filldraw[fill=white,draw=gray,thick] (#1,0) circle[radius=2.2pt];
    \node[below=4pt] at (#1,0) {#2};
  }
  \newcommand{\closedot}[3][magenta]{%
    \filldraw[fill=#1,draw=#1,thick] (#2,0) circle[radius=2.2pt];
    \node[below=4pt] at (#2,0) {#3};
  }
\usepackage[brazilian]{babel}
\usepackage{diagbox}
\usepackage{booktabs}

% collor-blind friendly pallete
\definecolor{figcolor1}{RGB}{0,70,70}
\definecolor{figcolor2}{RGB}{116,233,233}
\definecolor{figcolor3}{RGB}{146,0,0}
\definecolor{color1}{RGB}{0,70,70}
\definecolor{color2}{RGB}{116,233,233}
\definecolor{color3}{RGB}{146,0,0}


% A counter, since TikZ is not clever enough (yet) to handle
% arbitrary angle systems.
\newcount\mycount
\usepackage[percent]{overpic}
\usepackage{xcolor}
\usepackage{adjustbox}
\usepackage{bbold}
\newcommand{\rest}{\upharpoonright}

% symbols
\newcommand{\F}{\mathcal{F}}
\newcommand{\C}{\mathcal{C}}
\newcommand{\Open}{\mathcal{O}}
\newcommand{\A}{\mathcal{A}}
\newcommand{\B}{\mathcal{B}}
\newcommand{\G}{\mathcal{G}}
\newcommand{\M}{\mathcal{M}}
\newcommand{\N}{\mathbb{N}}
\newcommand{\Q}{\mathbb{Q}}
\newcommand{\meta}[1]{\overset{\text{meta}}{#1}}
\newcommand{\val}[1]{\llbracket #1 \rrbracket}
\newcommand{\valin}[3]{
	\sup_{#3 \in dom(#2)}\val{#3 = #1}#2(#3)}
\newcommand{\valsub}[3]{
	\inf_{#3 \in dom(#1)}(#1(#3)\implies \val{#3 \in #2})}
\newcommand{\valeq}[2]{\val{#1 \subset #2}\val{#2 \subset #1}}
% epimorphism arrow
\newcommand{\rightarrowdbl}{\rightarrow\mathrel{\mkern-14mu}\rightarrow}
\newcommand{\epimorphism}[2][]{%
	\xrightarrow[#1]{#2}\mathrel{\mkern-14mu}\rightarrow
}
\newcommand{\catname}[1]{\textnormal{\textsc{#1}}}

\newcommand{\R}{\mathbb{R}}
\newcommand{\D}{\mathcal{D}}
\newcommand{\Part}{\mathbb{P}}
\newcommand{\sequence}[1]{\langle #1 \rangle}
\DeclareMathOperator{\inter}{int}
\newcommand{\interior}[2]{\inter_{#2} (#1)}
\DeclareMathOperator{\clos}{cl}
\DeclareMathOperator{\diam}{diam}
\newcommand{\closure}[2]{\clos_{#2} (#1)}
\newcommand{\st}{\,:\,}
\newcommand{\dem}{\emph{dem.:}\\}
\renewcommand{\d}{\mathrm{d}}

% Theorem environments
\numberwithin{equation}{section}
% Theorem environments
\newtheorem{teor}{Teorema}[section]
\newtheorem{lema}{Lema}[section]
\newtheorem{prop}{Proposição}[section]
\newtheorem{properties}{Propriedades}[section]
\newtheorem{cor}{Corolário}[section]  
\newtheorem{defin}{Definição}[section]
\newtheorem{ex}{Exemplo}[section]
\newtheorem*{theorem*}{Teorema}
\newtheorem*{proposition*}{Proposição}
\newtheorem*{corollary*}{Corolário}
\newtheorem*{lemma*}{Lema}
\theoremstyle{remark}
\newtheorem{rmk}{Observação}[section]
\newtheorem*{question}{Problema}
\newtheorem{exerc}{Exercício}